%% kapa: exam=1 round=29 year=2018 period=2 form=A subject=accounting subject_ko=회계학 seq=6 number=46
%% source: Q-Net 감정평가사 기출문제 (cst00309) artlSeq=5202808 file=extracted/1008/5202808/2018년 제29회 감정평가사 1차시험 기출문제/2018년 제29회 감정평가사 1차시험 기출문제/2교시 A형.pdf
%% conversion: kordoc 4.24.51 -> tools/gen1.py (auto); review=proofread-sonnet; flags=none
\begin{question}{46}
재무보고를 위한 개념체계에서 재무제표 요소에 관한 설명으로 옳지 않은 것은?
\begin{choices}
  \item 자산이 갖는 미래경제적효익은 대체적인 제조과정의 도입으로 생산원가가 절감되는 경우와 같이 현금유출을 감소시키는 능력일 수도 있다.
  \item 자산의 존재를 판단하기 위해서 물리적 형태가 필수적인 것은 아니다.
  \item 경제적 효익에 대한 통제력은 법률적 권리의 결과이므로 법률적 통제가 있어야 자산의 정의를 충족시킬 수 있다.
  \item 기업은 일반적으로 구매나 생산을 통하여 자산을 획득하지만 다른 거래나 사건도 자산을 창출할 수 있다.
  \item 보증기간이 명백히 경과한 후에 발생하는 제품하자에 대해서도 수리해 주기로 방침을 정한 경우에 이미 판매된 제품과 관련하여 지출될 것으로 예상되는 금액은 부채이다.
\end{choices}
\end{question}
\answer{3}
